\section{خلاصة الفصل}

يُعدّ هذا الفصل حجر الأساس في الدراسة النظرية للذراع الروبوتية، إذ تناول النمذجة الحركية بوصفها الإطار الرياضي الذي يربط بين البنية الميكانيكية للذراع وبين سلوكها الحركي في فضاء العمل. وتنبع أهمية هذا الفصل من كونه يمهّد بصورة مباشرة للمراحل اللاحقة من الدراسة، سواء تلك المتعلقة بالتحليل التفاضلي، أو النمذجة الديناميكية، أو تصميم خوارزميات التحكم، أو حتى تطبيقات المحاكاة والتنفيذ العملي. فقبل الانتقال إلى أي مستوى متقدم في تحليل الذراع الروبوتية، لا بد أولًا من بناء وصف حركي دقيق يمكّن من فهم كيفية تموضع الوصلات والمفاصل، وكيفية انتقال الحركة من القاعدة حتى النهاية الفعالة.

وقد بدأ الفصل بتقديم تمهيد نظري يوضح مفهوم الحركة في الأنظمة الروبوتية، مع إبراز أن الذراع الروبوتية ليست مجرد مجموعة وصلات مترابطة ميكانيكيًا، بل هي منظومة هندسية يمكن توصيفها رياضيًا وفق نموذج منظم يسمح بتحليلها والتنبؤ بسلوكها. ومن هذا المنطلق، تم التركيز على توصيف الذراع الروبوتية من حيث عدد درجات الحرية، وطبيعة المفاصل المستخدمة، والعلاقات الهندسية التي تربط بين الأجزاء المختلفة. وقد أتاح هذا التوصيف الأولي تكوين صورة واضحة عن بنية الذراع، كما ساعد على تحديد المتغيرات الأساسية التي تحكم موضع النهاية الفعالة واتجاهها.

وفي هذا السياق، ظهرت الحاجة إلى اعتماد طريقة معيارية لوصف الذراع الروبوتية بصورة رياضية دقيقة ومنهجية، فتم استخدام اصطلاح \LR{Denavit-Hartenberg (DH)} بوصفه إحدى أكثر الطرق شيوعًا واعتمادًا في مجال الروبوتات. وتتمثل أهمية هذه الطريقة في قدرتها على اختزال الوصف الهندسي المعقد للذراع إلى مجموعة من البارامترات المحددة لكل وصلة أو مفصل، بحيث يصبح بالإمكان توصيف البنية الكاملة للذراع باستخدام جدول منظم يضم هذه المعاملات. ومن خلال هذا الأسلوب، تم تثبيت الإطارات المرجعية المحلية على المفاصل والوصلات بطريقة تراعي العلاقات الهندسية المتبادلة فيما بينها، الأمر الذي أتاح بناء جدول \LR{DH} الخاص بالذراع المدروسة.

ويُعتبر جدول \LR{DH} النتيجة المحورية الأولى في هذا الفصل، لأنه يمثل المرجع الذي انطلقت منه بقية عمليات الاشتقاق والتحليل. فمن خلال هذا الجدول، أمكن التعبير عن الانتقال من إطار مرجعي إلى آخر باستخدام مصفوفات التحويل المتجانس، وهي المصفوفات التي تجمع في آن واحد بين الدوران والإزاحة. ويُعد هذا الانتقال من الوصف الهندسي المباشر إلى الصياغة المصفوفية خطوة أساسية في بناء النموذج الرياضي للذراع، لأنه يجعل العلاقات بين الأجزاء المختلفة أكثر انتظامًا وقابلية للتحليل الجبري والبرمجي.

وانطلاقًا من ذلك، تناول الفصل اشتقاق مصفوفات التحويل الخاصة بكل مفصل على حدة، ثم تركيب هذه المصفوفات بالتسلسل للحصول على مصفوفة التحويل الكلية التي تصف وضعية النهاية الفعالة بالنسبة إلى الإطار المرجعي الأساسي. وهنا تجلّت دراسة الحركة الأمامية بوصفها العملية التي يتم فيها الانتقال من فضاء المفاصل إلى فضاء العمل، أي من قيم الزوايا أو الإزاحات المفصلية إلى الموضع والاتجاه النهائيين للنهاية الفعالة. وتُعد هذه الدراسة من أهم مسائل النمذجة الحركية، لأنها تسمح بتحديد موضع الذراع واتجاهها في أي لحظة تكون فيها قيم المتغيرات المفصلية معلومة.

وقد بينت نتائج الحركة الأمامية أن مصفوفة التحويل النهائية لا تقتصر فقط على تحديد إحداثيات الموضع، وإنما تحتوي كذلك على معلومات كاملة حول اتجاه النهاية الفعالة بالنسبة إلى القاعدة. وهذه النتيجة لها أهمية كبيرة من الناحية التطبيقية، لأن العديد من المهام الروبوتية لا تعتمد على الوصول إلى نقطة معينة فقط، وإنما تتطلب كذلك تحقيق اتجاه محدد للأداة أو الماسك النهائي. ومن ثم فإن دراسة الحركة الأمامية لا تمثل مجرد اشتقاق رياضي، بل تمثل أداة عملية أساسية في توصيف سلوك الذراع وتقييم قدرتها على تنفيذ المهام المطلوبة ضمن فضاء العمل.

كما تؤكد هذه الدراسة أن النموذج الحركي الأمامي هو الوسيلة المباشرة للتحقق من صحة البنية الرياضية المعتمدة، إذ يمكن من خلاله مقارنة النتائج النظرية مع التمركز الفيزيائي الفعلي للذراع أو مع نتائج المحاكاة الحاسوبية. ولذلك فإن نجاح اشتقاق الحركة الأمامية بصورة صحيحة يعد مؤشرًا على سلامة اختيار الإطارات المرجعية، وصحة جدول \LR{DH}، ودقة العلاقات الرياضية المستخدمة في تمثيل السلسلة الحركية.

وبعد تأسيس هذا الجانب، انتقل الفصل إلى دراسة الحركة العكسية، وهي من أكثر المسائل أهمية وتعقيدًا في مجال الروبوتات. فإذا كانت الحركة الأمامية تجيب عن سؤال: ما هو موضع واتجاه النهاية الفعالة عند معرفة متغيرات المفاصل؟ فإن الحركة العكسية تجيب عن السؤال المعاكس: ما هي قيم متغيرات المفاصل اللازمة لتحقيق موضع واتجاه مطلوبين مسبقًا للنهاية الفعالة؟ وتكمن صعوبة هذه المسألة في أن العلاقات الرياضية فيها غالبًا ما تكون غير خطية، وقد تؤدي إلى تعدد الحلول أو انعدامها أحيانًا ضمن المجال الفيزيائي المتاح للذراع.

وقد أظهر التحليل أن مسألة الحركة العكسية لا تتعلق فقط بإيجاد حل رياضي مجرد، وإنما تستوجب أيضًا مراعاة عدد من الجوانب العملية، مثل حدود حركة المفاصل، والقيود الميكانيكية، وتجنب الأوضاع غير الملائمة هندسيًا، إضافة إلى ضرورة اختيار الحل الأنسب من بين عدة حلول ممكنة. وهذا يعني أن دراسة الحركة العكسية تمثل مرحلة الربط الحقيقي بين النموذج الرياضي من جهة، والتنفيذ العملي للحركة من جهة أخرى. فالحل الذي يكون صحيحًا رياضيًا قد لا يكون بالضرورة هو الحل الأفضل من الناحية التطبيقية، إذا كان يؤدي مثلًا إلى تجاوز حدود أحد المفاصل، أو إلى اتخاذ وضعية غير آمنة، أو إلى مسار غير مناسب أثناء الحركة.

ومن النتائج المهمة التي أبرزها هذا الفصل أن الذراع الروبوتية قد تمتلك أكثر من تشكيل هندسي يحقق نفس موضع النهاية الفعالة، وهو ما يُعرف بتعدد حلول الحركة العكسية. وهذه الخاصية تُظهر الطبيعة غير الأحادية للعلاقة بين فضاء العمل وفضاء المفاصل، وتؤكد أن التحكم العملي في الذراع يتطلب وجود معايير إضافية لاختيار الحل الأنسب، مثل القرب من وضعية ابتدائية معينة، أو تقليل استهلاك الطاقة، أو تحسين انسيابية الحركة، أو الابتعاد عن الحالات الحرجة والمفاصل القريبة من حدودها.

ومن خلال الجمع بين نتائج الحركة الأمامية والحركة العكسية، يتضح أن الفصل قد أسّس بشكل متكامل للعلاقة الثنائية بين فضاء المفاصل وفضاء العمل. فالحركة الأمامية تمثل الاتجاه المباشر الذي يسمح بحساب موضع واتجاه النهاية الفعالة انطلاقًا من المتغيرات المفصلية، في حين تمثل الحركة العكسية الاتجاه المعاكس الذي يتيح تحديد القيم المفصلية اللازمة لتحقيق هدف حركي معين. وهذا التكامل يشكل جوهر النمذجة الحركية للذراع الروبوتية، ويُعد القاعدة التي تُبنى عليها تطبيقات التخطيط الحركي، وتوليد المسارات، والتحكم الموضعي، والتحكم في السرعة، والتحكم الذكي لاحقًا.

كذلك يمكن القول إن هذا الفصل لم يقتصر على عرض خطوات حسابية متسلسلة، بل قدّم تصورًا علميًا متكاملًا حول كيفية تحويل الذراع الروبوتية من منظومة ميكانيكية محسوسة إلى نموذج رياضي قابل للتحليل والمحاكاة والتنفيذ. وهذه النقلة من الوصف الفيزيائي إلى الصياغة الرياضية هي ما يمنح الدراسة طابعها الهندسي الحقيقي، لأن فهم النظام لم يعد مرتبطًا فقط بمكوناته المادية، بل أصبح قائمًا على العلاقات الرياضية التي تحكم عمله.

كما أن البناء المنهجي الذي تم اتباعه في هذا الفصل، بدءًا من المقدمة النظرية، ثم توصيف الذراع، فبناء جدول \LR{DH}، ثم اشتقاق الحركة الأمامية، وصولًا إلى تحليل الحركة العكسية، يعكس تسلسلًا علميًا منطقيًا ينسجم مع طبيعة الأنظمة الروبوتية ذات السلاسل الحركية. ويمنح هذا التسلسل القارئ فهمًا تراكميًا واضحًا، بحيث تكون كل خطوة ممهّدة لما يليها، وتكون النتائج النهائية ناتجة بصورة طبيعية عن المراحل السابقة.

ومن الناحية الأكاديمية، فإن هذا الفصل يمثل القاعدة المرجعية التي يمكن العودة إليها في جميع الفصول اللاحقة، لأن أي معالجة متقدمة للذراع الروبوتية ستعتمد، بشكل أو بآخر، على النموذج الحركي الذي تم تأسيسه هنا. فدراسة الجاكوبين مثلًا تنطلق من العلاقات الحركية بين المفاصل والنهاية الفعالة، ودراسة التفردات ترتبط ببنية النموذج الحركي نفسه، أما النمذجة الديناميكية فتفترض أساسًا أن البنية الحركية والهندسية للذراع قد تم تحديدها مسبقًا بصورة دقيقة. وكذلك الأمر بالنسبة لخوارزميات التحكم، إذ لا يمكن تصميم متحكم فعّال دون معرفة واضحة للعلاقة بين أوامر المفاصل والوضعية الناتجة للنهاية الفعالة.

وبناءً على جميع ما سبق، يمكن الاستنتاج أن هذا الفصل قد نجح في بناء إطار نظري ورياضي متكامل للنمذجة الحركية للذراع الروبوتية. فقد تم فيه الانتقال بصورة منهجية من التوصيف البنيوي للذراع إلى التمثيل الرياضي للعلاقات بين أجزائها، ثم إلى اشتقاق المعادلات التي تصف حركتها المباشرة والعكسية. وبهذا المعنى، فإن ما تم التوصل إليه في هذا الفصل لا يمثل فقط خلاصة للدراسة الحركية، بل يشكل أيضًا قاعدة تأسيسية ضرورية لكل ما سيلي من دراسات تحليلية أو تطبيقية تخص الذراع الروبوتية.

وعليه، فإن هذا الفصل يرسّخ الفهم العلمي للعلاقة بين التركيب الميكانيكي للذراع وبين أدائها الحركي، ويبرز أن النمذجة الحركية ليست مجرد مرحلة تمهيدية، بل هي جوهر أساسي في دراسة الأنظمة الروبوتية، وعنصر لا غنى عنه في أي محاولة جادة لفهم الذراع الروبوتية أو تطويرها أو التحكم بها بكفاءة ودقة.